% Work-MAT | 2. rocnik | pripravovana hodina 30
% Návaznost: uživatelské připravy.zip (hodiny 01–20), autorská učebnice na GitHubu.
% Realisticky.cz: didaktická inspirace v posloupnosti a typech úloh; zadání jsou vlastní.
\documentclass[11pt,a4paper]{article}
\input{../../sablony/pripravy/priprava.sty}
\nastavPripravu{2. ročník}{Exponenciální nerovnice}{30}{Exponenciální nerovnice II}
\begin{document}
\blok{Motivace a návaznost}
Navážeme na monotonii a přidáme nerovnice, v nichž nejprve upravujeme mocniny na společný základ. Po výpočtu řešení zapíšeme intervalem.
\poznamka{Orientační nabídka pro 45--55 minut: krátké zopakování a motivace, nový výklad, tři společné vzorové úlohy, samostatná práce. Vyber podle tempa třídy; zbývající zadání slouží jako rezerva.}
\blok{Výklad}
Při kladném základu $a>1$ se nerovnost exponentů neobrací; u $0<a<1$ se obrací. Ekvivalentní úpravy exponentů jsou stejné jako u rovnic.
\blok{Společné příklady}
\begin{prikladReseniBox}[title={Příklad 1}]
Vyřeš, resp. urči: $8^{x-1}\le4^{x+1}$.
\tcblower
\[\begin{aligned}2^{3x-3}\le2^{2x+2} \\ 3x-3\le2x+2 \\ x\le5\end{aligned}\]
\end{prikladReseniBox}
\begin{prikladReseniBox}[title={Příklad 2}]
Vyřeš, resp. urči: $\left(\frac{1}{2}\right)^{2x-3}>\frac{1}{8}$.
\tcblower
\[\begin{aligned}\left(\frac{1}{2}\right)^{2x-3}>\left(\frac{1}{2}\right)^3 \\ 2x-3<3 \\ x<3\end{aligned}\]
\end{prikladReseniBox}
\poznamka{Při úlohách se stejnými výrazy v exponentu upozorni na nerovnosti pravdivé pro všechna $x$ nebo pro žádné $x$.}
\blok{Samostatné procvičování}
\begin{zadaniReseniBox}[title={Samostatná práce 1}]
\begin{enumerate}[label=\alph*)]
\item $9^x>3^{x+3}$.
\item $8^{x-1}\le4^{x+1}$.
\item $(\frac{1}{2})^{2x-3}>(\frac{1}{2})^3$.
\item $4^{x+2}<2^{3x}$.
\item $25^x\ge5^{x+3}$.
\item $(\frac{1}{3})^{x-1}\le(\frac{1}{3})^2$.
\end{enumerate}
\tcblower
\begin{enumerate}[label=\alph*)]
\item $x>3$.
\item $x\le5$.
\item $x<3$.
\item $x>4$.
\item $x\ge3$.
\item $x\ge3$.
\end{enumerate}
\end{zadaniReseniBox}
\blok{Dobrovolné rozšíření}
\blok{Úkol na doma}
Při početních úlohách napiš postup, zvolenou metodu a zkoušku v původním zadání. U soustav kontroluj obě rovnice; u nerovnic výsledek zakresli na osu.
\begin{zadaniReseniBox}[title={Domácí úkol}]
\begin{enumerate}[label=\alph*)]
\item $4^x\ge2^{x+2}$.
\item $8^{x-1}\le4^{x+1}$.
\item $9^{x-1}>3^{x+2}$.
\item $(\frac{1}{2})^{2x}\le(\frac{1}{2})^4$.
\item $16^x>8^{x+1}$.
\item $5^{2x+1}<25^{x+1}$.
\item Vysvětli poslední případ..
\end{enumerate}
\tcblower
\begin{enumerate}[label=\alph*)]
\item $x\ge2$.
\item $x\le5$.
\item $x>4$.
\item $x\ge2$.
\item $x>3$.
\item všechna x.
\item $2x+1<2x+2$ platí vždy..
\end{enumerate}
\end{zadaniReseniBox}
\poznamka{Domácí úkol obsahuje 7 položek; předpokládaná domácí práce přibližně 15--30 minut. Po větších blocích jsou vědomě promíchané dříve probrané typy.}
\end{document}
